\documentclass[10pt,a4paper]{amsart}
\usepackage[margin=19mm]{geometry}
\usepackage{amsmath,amssymb,mathtools}
\usepackage[scr=rsfs,cal=euler]{mathalfa}
\usepackage{xfrac}
\usepackage[unicode,hidelinks,bookmarks=false]{hyperref}
\newcommand{\ie}{i.e.\ }
\newcommand{\cf}{cf.\ }
\newcommand{\resp}{respectively\ }
\newcommand{\Subsection}{\subsection}
\providecommand{\nobreakdash}{\nobreak\hbox{-}}
\setlength{\emergencystretch}{2em}
\multlinegap0pt
\usepackage{bm}
\usepackage{bbm}
\usepackage{dsfont}
\usepackage{xspace}
\usepackage{cancel}
\usepackage{mathtools}
\usepackage{amsthm}
\makeatletter
\@ifundefined{theorem}{\newtheorem{theorem}{Theorem}}{}
\@ifundefined{lemma}{\newtheorem{lemma}[theorem]{Lemma}}{}
\makeatother
\newcommand{\parenthesis}[1]{\left(#1\right)} % round bracket
\renewcommand{\to}{\rightarrow}
\title[Classification of singular limits for free boundary and sing]{Classification of singular limits for free boundary and singularly perturbed elliptic problems: the Dancer-Yan spikes revisited}
\author{Daniele Bartolucci, Aleks Jevnikar, Juncheng Wei, Ruijun Wu}
\date{Source: arXiv:2507.20725v1 (CC BY 4.0)}
\begin{document}
\maketitle

\noindent\emph{Source and licence.} Classification of singular limits for free boundary and singularly perturbed elliptic problems: the Dancer-Yan spikes revisited. Version arXiv:2507.20725v1 (CC BY 4.0), distributed under the Creative Commons Attribution 4.0 International licence, as recorded in the arXiv licence metadata for this exact version and shown on the abstract page https://arxiv.org/abs/2507.20725v1. Full rights notice: licenses/arxiv-2507.20725-CC-BY-4.0.txt.\par\smallskip

\noindent\textit{Editorial scope.} Counted unit: the Lemma whose source label is \texttt{lemma:rel dist} in the article (source0.tex lines 1757--1762, source label \texttt{lemma:rel dist}), delivered together with the article's own complete original proof of it (source0.tex lines 1764--1786, one \texttt{proof} environment of 23 lines). No mathematics is written, repaired or re-derived: statement and proof are the article's own text line for line. Required context retained uncounted: eq:un def (lines 579--581); eq:un2 def (lines 1727--1729). Every definition, display and label of those lines is retained unchanged. Omitted from this package and not redistributed: 1--578, 582--1726, 1730--1756, 1763--1763, 1787--2447. Nothing inside a retained excerpt is omitted or altered: every retained body is delivered exactly as the article's lines are written, so no omission receipt of any kind accompanies this package. Citations: the retained excerpts carry no citation command at all, so no bibliography entry of the article is transcribed and no reference is redistributed. Reference adaptation: none was needed and none was made; every reference inside the delivered unit resolves inside it, so no unresolved reference or citation is produced. Printed slips kept literal: no printed word, symbol, hypothesis or number of the counted statement or of its proof was altered. Layout and class: the article's own class is \texttt{amsart} with packages including color, hyperref, scalerel, amsmath, amssymb, amsthm, amsfonts, mathrsfs, amsopn, mathtools, inputenc, srcltx, graphicx, exscale; the wrapper is the lane's pinned \texttt{amsart} editorial wrapper at 10pt on A4, which restates the theorem-like environments the retained text needs and adds the article's own macro definitions that the retained text needs (\texttt{\textbackslash parenthesis}, \texttt{\textbackslash to}), with extra wrapper packages (bm, bbm, dsfont, xspace, cancel, mathtools). The delivered numbering is the unit's own. Assets: no retained span contains a figure environment, a graphics-inclusion command, a PDF-page inclusion, a table or a TikZ picture; no image file is copied into this package, the package declares no asset dependency and no image file is redistributed with it. Licence: version arXiv:2507.20725v1 of this article is distributed under the Creative Commons Attribution 4.0 International licence, as recorded in the arXiv licence metadata for this exact version and shown on the abstract page https://arxiv.org/abs/2507.20725v1. A full rights notice is delivered at licenses/arxiv-2507.20725-CC-BY-4.0.txt. Characters: every retained excerpt is pure ASCII, so the delivered engine, XeLaTeX with its default Latin Modern text font, reports no missing character.
            \begin{align}\label{eq:un def}
                u_n(z)= {\mathcal{R}}_{t_n}\tilde{v}_n(z):= t_n^{\frac{2}{p-1}}\tilde{v}_n(t_n z)
            \end{align}

    \begin{align}\label{eq:un2 def}
        u_{n,2}(z)\coloneqq t_{n,2}^{\frac{2}{p-1}} \theta_n \parenthesis{ v_n(x_{n,2}+ s_{n} t_{n,2} z)-1 }, \qquad z\in B_{{2\delta} / {s_n t_{n,2}}}(0).
    \end{align}

    \begin{lemma}\label{lemma:rel dist}
       Let $x_{n,2}$, $x_{n,1}$ and $s_n t_{n,2}$ be defined as above, then we have,
        \begin{align}\label{eq:far I}
            \frac{|x_{n,1}-x_{n,2}|}{s_n t_{n,2}}\to+\infty.
        \end{align}
    \end{lemma}

    \begin{proof}
        By contradiction assume that there exists~$C>0$ such that,
        \begin{align*}
            \frac{R_{n,1} t_{n,1}}{t_{n,2}}\leq \frac{|x_{n,1}-x_{n,2}|}{s_n t_{n,2}}\leq C,
        \end{align*}
        implying in particular that,
        \begin{align*}
          2R_{n,1}\leq \frac{|x_{n,1}-x_{n,2}|}{s_n} \leq C t_{n,2},
        \end{align*}
        and $\frac{|x_{n,1}-x_{n,2}|}{s_n}\to +\infty$ and~$t_{n,2}\to +\infty$ as~$n\to +\infty$.

        Recall the functions~$u_{n,1}$ and~$u_{n,2}$ defined by~\eqref{eq:un def} and~\eqref{eq:un2 def}.
        Putting
        \begin{align*}
            z_{n,1,2}\coloneqq \frac{x_{n,1} -x_{n,2}}{s_n t_{n,2}},
        \end{align*}
        which is uniformly bounded, possibly along a subsequence we have $z_{n,1, 2}\to z_{\infty,1,2}$, and then
        \begin{align*}
            u_{n,2}(z_{n,1,2})= t_{n,2}^{\frac{2}{p-1}}\theta_n\parenthesis{ v_n(x_{n,1} )-1}= \parenthesis{\frac{t_{n,2}}{t_{n,1}}}^{\frac{2}{p-1}} u_{n,1}(0)\to +\infty
        \end{align*}
        since~$u_{n,1}(0)=\phi(0)>0$ and~$\frac{t_{n,2}}{t_{n,1}}\to +\infty$.
        This is a contradiction to the fact that~$u_{n,2}$ is bounded from above by~$\phi(0)$.
    \end{proof}

\end{document}
